\chapter{Technische Analyse}

Im Mittelpunkt der Anforderungen steht, dass die Häuser entsprechend des Region Scores optimal gruppiert werden sollen(REQ-xx). Dabei stehen zwei Themenbereiche hervor. DIe MAthemathische Berechnungsgrundlage des Region Scors und die Gruppierungsalgorithmen selbst.

\subsection{Region Score}
\subsubsection{Autarky}
\subsubsection{}

\subsection{Gruppierungsalgorithmen}
ES sollen verschiedene herangehensweisen zum Gruppieren der Häuser zu Energiecommunities verglichen werden. Es gibt mehrere Ansätze in der Informatik Datenpunkte zu gruppieren. In dieser Arbeit werden folgende Algorithmen miteinander verglichen:

Eine Brute Force Methode, bei der alle möglichen Gruppen evaluiert werden.
2.K-Means als Vertreter einers klassischen Clusterverfahren.
Hierarchisch Bottom up und Top down
Genetisches verfahren als vertreter von Zufallsmethoden

\subsubsection{Brute Force}
\subsubsection{K-Means}
\subsubsection{Hierarchisch}
\subsubsection{Genetisch}


